% ============================================================
%  Chapitre 4 — Architecture et Conception
% ============================================================
\chapter{Architecture et Conception du Système}
\thispagestyle{fancy}

\section{Introduction}

Ce chapitre détaille l'architecture logicielle et matérielle de \farmAI{} v3.0.
L'approche de conception repose sur trois principes directeurs :

\begin{enumerate}
  \item \textbf{Séparation des responsabilités (SoC)} — chaque couche a une
        responsabilité unique et bien délimitée.
  \item \textbf{Principe de dépendances inversées (DIP)} — les couches hautes ne
        dépendent pas des implémentations concrètes des couches basses.
  \item \textbf{Multi-tenancy par isolation} — chaque exploitation agricole est isolée
        par un \code{farm\_id} propagé dans toutes les requêtes.
\end{enumerate}

\section{Diagrammes de cas d'utilisation}

\subsection{Cas d'utilisation global}

La figure \ref{fig:uc_global} présente le diagramme de cas d'utilisation global de
\farmAI{} v3.0, avec ses quatre acteurs principaux et les 15 cas d'utilisation
couverts par la plateforme.

\begin{figure}[H]
\centering
\resizebox{\textwidth}{!}{%
% ============================================================
%  Diagramme de Cas d'Utilisation Global — Smart Farm AI v3.0
%  Acteurs : Owner, Worker, ESP32 Nodes, Ollama LLM
% ============================================================
\begin{tikzpicture}[
  x=1cm, y=1cm,
  actor/.style={
    font=\scriptsize\bfseries, align=center,
    text width=1.8cm
  },
  usecase/.style={
    ellipse, draw=farmblue!70, fill=farmblue!8,
    text width=2.5cm, align=center,
    font=\tiny, inner sep=2pt,
    minimum height=0.7cm
  },
  ucsmall/.style={
    ellipse, draw=farmgreen!70, fill=farmgreen!8,
    text width=2.2cm, align=center,
    font=\tiny, inner sep=2pt,
    minimum height=0.6cm
  },
  include/.style={dashed, ->, >=Stealth, farmblue!60, font=\tiny},
  extend/.style={dashed, ->, >=Stealth, farmorange!70, font=\tiny},
  assoc/.style={-, thick, farmgray!60}
]

% ── Frontière système ────────────────────────────────────────
\draw[rounded corners=12pt, draw=farmblue, fill=farmblue!3, line width=1.5pt]
  (1.2, -9.5) rectangle (12.8, 1.2);
\node[font=\normalsize\bfseries\color{farmblue}, above] at (7, 1.2)
  {\farmAI{} v3.0 — Système};

% ── Macro stickman TikZ ──────────────────────────────────────
\newcommand{\stickman}[3]{% x, y, label
  \begin{scope}[shift={(#1,#2)}]
    \draw[thick, farmgray] (0,0.55) circle (0.18);
    \draw[thick, farmgray] (0,0.37) -- (0,-0.15);
    \draw[thick, farmgray] (-0.22,0.15) -- (0.22,0.15);
    \draw[thick, farmgray] (0,-0.15) -- (-0.18,-0.55);
    \draw[thick, farmgray] (0,-0.15) -- ( 0.18,-0.55);
  \end{scope}
  \node[actor] at (#1, #2-0.85) {#3};
}

% ── Acteur : Owner (gauche haut) ─────────────────────────────
\stickman{-1.5}{-0.7}{Owner\\(Propriétaire)}
\coordinate (owner) at (-1.5, -1.5);

% ── Acteur : Worker (gauche bas) ─────────────────────────────
\stickman{-1.5}{-6.7}{Worker\\(Ouvrier)}
\coordinate (worker) at (-1.5, -7.5);

% ── Acteur : ESP32 Nodes (droite haut) ───────────────────────
\stickman{15.5}{-0.7}{ESP32 Nodes\\(NODE\_A/B)}
\coordinate (esp32) at (15.5, -1.5);

% ── Acteur : Ollama LLM (droite bas) ─────────────────────────
\stickman{15.5}{-6.7}{Ollama LLM\\(Labess-7B)}
\coordinate (ollama) at (15.5, -7.5);

% ── Cas d'utilisation Owner ───────────────────────────────────
\node[usecase] (uc_auth)    at (3.5, 0)    {S'authentifier\\(JWT + OTP)};
\node[usecase] (uc_farm)    at (3.5, -1.8) {Gérer les\\Fermes};
\node[usecase] (uc_dash)    at (3.5, -3.5) {Tableau de bord\\(KPIs + Météo)};
\node[usecase] (uc_tele)    at (3.5, -5.2) {Visualiser\\Télémétrie IoT};
\node[usecase] (uc_alert)   at (3.5, -6.8) {Recevoir\\Alertes IoT};
\node[usecase] (uc_map)     at (3.5, -8.5) {Consulter\\Map SIG};

% ── Cas d'utilisation partagés Owner/Worker ────────────────────
\node[usecase] (uc_scan)    at (7,   -1.5) {Scanner\\(Vision YOLO)};
\node[usecase] (uc_hive)    at (7,   -4.5) {Gestion\\Apiculture};
\node[usecase] (uc_worker)  at (7,   -7.5) {Gestion\\Ouvriers};

% ── Cas d'utilisation Assistant ───────────────────────────────
\node[usecase] (uc_chat)    at (10.5, -1.5) {Dialoguer\\(Darija)};
\node[ucsmall] (uc_stt)     at (10.5, -3.0) {Parler / Écouter\\(STT/TTS)};
\node[ucsmall] (uc_vision)  at (10.5, -4.5) {Analyse Image\\(LLaVA)};

% ── Cas d'utilisation IoT / System ────────────────────────────
\node[usecase] (uc_iot_post) at (10.5, -6.5) {Envoyer\\Télémétrie};
\node[ucsmall] (uc_anomaly)  at (10.5, -8.5) {Détecter\\Anomalies};

% ── Associations Owner ────────────────────────────────────────
\draw[assoc] (owner) -- (uc_auth);
\draw[assoc] (owner) -- (uc_farm);
\draw[assoc] (owner) -- (uc_dash);
\draw[assoc] (owner) -- (uc_tele);
\draw[assoc] (owner) -- (uc_alert);
\draw[assoc] (owner) -- (uc_map);
\draw[assoc] (owner) -- (uc_scan);
\draw[assoc] (owner) -- (uc_hive);
\draw[assoc] (owner) -- (uc_worker);
\draw[assoc] (owner) -- (uc_chat);

% ── Associations Worker ───────────────────────────────────────
\draw[assoc] (worker) -- (uc_auth);
\draw[assoc] (worker) -- (uc_scan);
\draw[assoc] (worker) -- (uc_worker);

% ── Associations ESP32 ────────────────────────────────────────
\draw[assoc] (esp32) -- (uc_iot_post);

% ── Associations Ollama ───────────────────────────────────────
\draw[assoc] (ollama) -- (uc_chat);

% ── Relations include / extend ────────────────────────────────
\draw[include] (uc_chat)    -- node[above, sloped]{\texttt{<<include>>}} (uc_stt);
\draw[extend]  (uc_chat)    -- node[right]{\texttt{<<extend>>}}          (uc_vision);
\draw[include] (uc_dash)    -- node[above, sloped, pos=0.4]{\texttt{<<include>>}} (uc_tele);
\draw[extend]  (uc_iot_post)-- node[right]{\texttt{<<extend>>}}          (uc_anomaly);

\end{tikzpicture}

}
\caption{Diagramme de cas d'utilisation global — Smart Farm AI v3.0}
\label{fig:uc_global}
\end{figure}

\begin{table}[H]
\centering
\caption{Description des cas d'utilisation principaux}
\label{tab:uc_desc}
\small
\begin{tabular}{cllp{5cm}}
\toprule
\textbf{UC} & \textbf{Acteur} & \textbf{Intitulé} & \textbf{Description} \\
\midrule
UC01 & Owner/Worker & S'authentifier &
  Connexion via email/mdp + JWT 30j, OTP email 6 chiffres \\
UC02 & Owner & Gérer les Fermes &
  CRUD fermes, associations GPS, dashboard multi-fermes \\
UC03 & Owner & Tableau de bord &
  KPIs temps réel : météo, IoT live, 6 badges espèces \\
UC04 & Owner & Visualiser Télémétrie &
  Courbes historiques Recharts, filtre NODE\_A/NODE\_B \\
UC05 & Owner & Recevoir Alertes &
  SWARM\_RISK, THERMAL\_FAULT via WebSocket \\
UC06 & Owner/Worker & Scanner (YOLO) &
  Inférence 13 modèles, caméra live ou upload image \\
UC07 & Owner & Gestion Apiculture &
  HiveIntel 7 sous-modules (voir UC apiculture) \\
UC08 & Owner & Gestion Ouvriers &
  CRUD ouvriers, assignation tâches, suivi GPS \\
UC09 & Owner & Consulter Map SIG &
  Leaflet, découverte vétérinaires/fournisseurs OSM \\
UC10 & Owner & Dialoguer Darija &
  Assistant Labess-7B, STT/TTS, vision multimodale \\
UC11 & ESP32 & Envoyer Télémétrie &
  HTTP POST JSON toutes les 10 secondes \\
UC12 & System & Détecter Anomalies &
  IsolationForest sur flux temps réel \\
\bottomrule
\end{tabular}
\end{table}

\subsection{Cas d'utilisation — Module Apiculture HiveIntel}

Le module apiculture constitue l'une des contributions les plus originales du projet,
intégrant la gestion complète d'un rucher avec la télémétrie IoT du NODE\_B.

\begin{figure}[H]
\centering
\resizebox{0.95\textwidth}{!}{%
% ============================================================
%  Diagramme de Cas d'Utilisation — Module Apiculture HiveIntel
%  Acteurs : Owner Apiculteur, NODE_B ESP32
% ============================================================
\begin{tikzpicture}[
  x=1cm, y=1cm,
  actor/.style={font=\scriptsize\bfseries, align=center, text width=1.8cm},
  usecase/.style={
    ellipse, draw=farmgreen!80!black, fill=farmgreen!10,
    text width=2.8cm, align=center,
    font=\tiny, inner sep=3pt, minimum height=0.75cm
  },
  uciot/.style={
    ellipse, draw=farmorange!80!black, fill=farmorange!10,
    text width=2.8cm, align=center,
    font=\tiny, inner sep=3pt, minimum height=0.75cm
  },
  include/.style={dashed, ->, >=Stealth, farmblue!60, font=\tiny},
  extend/.style={dashed, ->, >=Stealth, farmorange!70, font=\tiny},
  assoc/.style={-, thick, farmgray!60}
]

% ── Frontière système ────────────────────────────────────────
\draw[rounded corners=10pt, draw=farmgreen, fill=farmgreen!3, line width=1.5pt]
  (1.5, -11.8) rectangle (13, 1.0);
\node[font=\large\bfseries\color{farmgreen}, above] at (7.25, 1.0)
  {HiveIntel — Module Apiculture};

% ── Acteur : Owner Apiculteur ─────────────────────────────────
\begin{scope}[shift={(-0.5,-4.8)}]
  \draw[thick,farmgray] (0,0.55) circle (0.18);
  \draw[thick,farmgray] (0,0.37)--(0,-0.15);
  \draw[thick,farmgray] (-0.22,0.15)--(0.22,0.15);
  \draw[thick,farmgray] (0,-0.15)--(-0.18,-0.55);
  \draw[thick,farmgray] (0,-0.15)--( 0.18,-0.55);
\end{scope}
\node[actor] (owner) at (-0.5, -5.9) {Owner\\Apiculteur};

% ── Acteur : NODE_B ESP32 ─────────────────────────────────────
\begin{scope}[shift={(15,-7.8)}]
  \draw[thick,farmgray] (0,0.55) circle (0.18);
  \draw[thick,farmgray] (0,0.37)--(0,-0.15);
  \draw[thick,farmgray] (-0.22,0.15)--(0.22,0.15);
  \draw[thick,farmgray] (0,-0.15)--(-0.18,-0.55);
  \draw[thick,farmgray] (0,-0.15)--( 0.18,-0.55);
\end{scope}
\node[actor] (nodeb) at (15, -9.2) {NODE\_B\\ESP32 Rucher};

% ── Cas d'utilisation Owner ───────────────────────────────────
\node[usecase] (uc_apiaire)   at (4.5, 0)    {Créer / Gérer\\Apiaire};
\node[usecase] (uc_ruche)     at (4.5, -1.9) {Ajouter Ruche\\+ QR Code};
\node[usecase] (uc_visite)    at (4.5, -3.8) {Enregistrer\\Visite Ruche};
\node[usecase] (uc_recolte)   at (4.5, -5.7) {Gérer\\Récoltes Miel};
\node[usecase] (uc_stock)     at (4.5, -7.6) {Gérer\\Stock Hausses};
\node[usecase] (uc_depense)   at (4.5, -9.5) {Enregistrer\\Dépenses};
\node[usecase] (uc_planning)  at (4.5, -11.2){Planifier\\Interventions};

% ── Cas d'utilisation Analytics ───────────────────────────────
\node[usecase] (uc_analytics) at (9.5, -1.9) {Analytics\\Apiculture};
\node[usecase] (uc_rapport)   at (9.5, -3.8) {Générer\\Rapport PDF};
\node[uciot]   (uc_tele)      at (9.5, -5.7) {Visualiser\\Télémétrie NODE\_B};
\node[uciot]   (uc_alerte)    at (9.5, -7.6) {Recevoir Alertes\\(SWARM/THERMAL)};

% ── Cas d'utilisation NODE_B ──────────────────────────────────
\node[uciot]   (uc_send)      at (9.5, -9.5)  {Envoyer\\Métriques IoT};
\node[uciot]   (uc_check)     at (9.5, -11.2) {Contrôler\\LEDs Alertes};

% ── Associations Owner ────────────────────────────────────────
\foreach \uc in {uc_apiaire, uc_ruche, uc_visite, uc_recolte,
                 uc_stock, uc_depense, uc_planning,
                 uc_analytics, uc_rapport, uc_tele, uc_alerte}
  \draw[assoc] (owner) -- (\uc);

% ── Associations NODE_B ───────────────────────────────────────
\draw[assoc] (nodeb) -- (uc_send);
\draw[assoc] (nodeb) -- (uc_check);

% ── Relations include ─────────────────────────────────────────
\draw[include] (uc_analytics) --
  node[above, sloped, font=\tiny]{\texttt{<<include>>}} (uc_tele);
\draw[include] (uc_send) --
  node[above, sloped, font=\tiny]{\texttt{<<include>>}} (uc_tele);

% ── Relations extend ──────────────────────────────────────────
\draw[extend] (uc_tele) --
  node[right, font=\tiny]{\texttt{<<extend>>}} (uc_alerte);
\draw[extend] (uc_visite) --
  node[above, sloped, font=\tiny]{\texttt{<<extend>>}} (uc_planning);
\draw[extend] (uc_analytics) --
  node[right, font=\tiny]{\texttt{<<extend>>}} (uc_rapport);

% ── Légende ───────────────────────────────────────────────────
\node[usecase, scale=0.75] at (3, -12.8)   {UC Owner};
\node[uciot,   scale=0.75] at (7.5, -12.8) {UC IoT};
\node[font=\tiny, right] at (4.4, -12.8) {Fonctions apicoles};
\node[font=\tiny, right] at (8.9, -12.8) {Télémétrie ESP32};

\end{tikzpicture}

}
\caption{Diagramme de cas d'utilisation — Module HiveIntel (Apiculture)}
\label{fig:uc_bee}
\end{figure}

\section{Architecture système en couches}

\farmAI{} v3.0 adopte une architecture en \textbf{5 couches} inspirée du modèle
N-tiers, avec une séparation nette entre la présentation, la logique métier, l'accès
aux données et l'infrastructure physique.

\begin{figure}[H]
\centering
\resizebox{\textwidth}{!}{%
% ============================================================
%  Architecture Système — Smart Farm AI v3.0 (5 couches)
% ============================================================
\begin{tikzpicture}[
  x=1cm, y=1cm,
  layer/.style={
    rectangle, rounded corners=6pt,
    draw=#1!70!black, fill=#1!15,
    text width=16cm, minimum height=1.6cm,
    align=center, inner sep=6pt,
    font=\small\bfseries
  },
  component/.style={
    rectangle, rounded corners=3pt,
    draw=#1!60!black, fill=#1!20,
    text width=3cm, minimum height=1.1cm,
    align=center, inner sep=4pt,
    font=\tiny\bfseries
  },
  arrow/.style={->, thick, >=Stealth, #1},
  label/.style={font=\scriptsize\bfseries\color{#1}, sloped}
]

% ─── Couche 1 : Clients ──────────────────────────────────────
\node[draw=farmblue, fill=farmblue!8, rounded corners=6pt,
      minimum width=16cm, minimum height=1.8cm,
      align=center, font=\small\bfseries\color{farmblue}]
  (clients) at (0, 0) {};
\node[above=-0.3cm of clients, font=\tiny\bfseries\color{farmblue!60}]
  {COUCHE CLIENTS};
\node[component=farmblue] at (-5.5,  0) {Owner Dashboard\\React/Vite SPA};
\node[component=farmblue] at (-1.8,  0) {Worker PWA\\Offline (Dexie)};
\node[component=farmblue] at ( 1.8,  0) {Map Center\\Leaflet + OSM};
\node[component=farmblue] at ( 5.5,  0) {Assistant\\Chat + STT/TTS};

% ─── Flèche Clients → API ────────────────────────────────────
\draw[arrow=farmblue] (0, -0.9) -- (0, -1.7)
  node[midway, right, font=\tiny] {HTTPS / REST / WebSocket};

% ─── Couche 2 : API Gateway ──────────────────────────────────
\node[draw=farmgreen, fill=farmgreen!8, rounded corners=6pt,
      minimum width=16cm, minimum height=1.8cm,
      align=center, font=\small\bfseries\color{farmgreen}]
  (api) at (0, -2.6) {};
\node[above=-0.3cm of api, font=\tiny\bfseries\color{farmgreen!60}]
  {COUCHE API GATEWAY};
\node[component=farmgreen] at (-5.5, -2.6) {JWT Middleware\\Auth + OTP};
\node[component=farmgreen] at (-1.8, -2.6) {FastAPI\\50+ Endpoints};
\node[component=farmgreen] at ( 1.8, -2.6) {WebSocket\\Alertes Temps Réel};
\node[component=farmgreen] at ( 5.5, -2.6) {CORS + Rate\\Limiting};

% ─── Flèche API → Services ───────────────────────────────────
\draw[arrow=farmgreen] (0, -3.5) -- (0, -4.3)
  node[midway, right, font=\tiny] {Appels internes Python};

% ─── Couche 3 : Services ─────────────────────────────────────
\node[draw=farmorange, fill=farmorange!8, rounded corners=6pt,
      minimum width=16cm, minimum height=1.8cm,
      align=center, font=\small\bfseries\color{farmorange}]
  (services) at (0, -5.2) {};
\node[above=-0.3cm of services, font=\tiny\bfseries\color{farmorange!60}]
  {COUCHE SERVICES MÉTIER};
\node[component=farmorange] at (-6.3, -5.2) {CV Service\\YOLO ×8};
\node[component=farmorange] at (-3.2, -5.2) {RAG Service\\ChromaDB};
\node[component=farmorange] at ( 0,   -5.2) {LLM Agent\\Labess-7B};
\node[component=farmorange] at ( 3.2, -5.2) {Anomaly\\IsolationForest};
\node[component=farmorange] at ( 6.3, -5.2) {Weather\\OpenMeteo};

% ─── Flèche Services → Données ───────────────────────────────
\draw[arrow=farmorange] (0, -6.1) -- (0, -6.9)
  node[midway, right, font=\tiny] {SQLAlchemy ORM / direct calls};

% ─── Couche 4 : Données ──────────────────────────────────────
\node[draw=farmred, fill=farmred!8, rounded corners=6pt,
      minimum width=16cm, minimum height=1.8cm,
      align=center, font=\small\bfseries\color{farmred}]
  (data) at (0, -7.8) {};
\node[above=-0.3cm of data, font=\tiny\bfseries\color{farmred!60}]
  {COUCHE DONNÉES \& MODÈLES};
\node[component=farmred] at (-6.3, -7.8) {PostgreSQL\\30+ Tables};
\node[component=farmred] at (-3.2, -7.8) {ChromaDB\\HNSW cosine};
\node[component=farmred] at ( 0,   -7.8) {MLflow\\Registry};
\node[component=farmred] at ( 3.2, -7.8) {DVC\\Data Versions};
\node[component=farmred] at ( 6.3, -7.8) {Ollama\\Models};

% ─── Flèche Données → IoT ────────────────────────────────────
\draw[arrow=farmred] (0, -8.7) -- (0, -9.5)
  node[midway, right, font=\tiny] {HTTP POST /api/v1/iot/telemetry};

% ─── Couche 5 : IoT ──────────────────────────────────────────
\node[draw=farmgray, fill=farmgray!10, rounded corners=6pt,
      minimum width=16cm, minimum height=1.8cm,
      align=center, font=\small\bfseries\color{farmgray}]
  (iot) at (0, -10.4) {};
\node[above=-0.3cm of iot, font=\tiny\bfseries\color{farmgray!80}]
  {COUCHE IOT (TERRAIN)};
\node[component=farmgray] at (-5.5, -10.4)
  {NODE\_A ESP32\\Irrigation\\(sol, pression, débit)};
\node[component=farmgray] at (-1.8, -10.4)
  {Wokwi Simulation\\WiFi Virtuel\\Wokwi-GUEST};
\node[component=farmgray] at ( 1.8, -10.4)
  {NODE\_B ESP32\\Rucher\\(poids, temp, hum)};
\node[component=farmgray] at ( 5.5, -10.4)
  {PlatformIO\\Build + Flash\\OTA};

% ─── Flèche IoT ← retour alertes ─────────────────────────────
\draw[arrow=farmblue!50, dashed, bend right=70]
  (6, -7.8) to node[right, font=\tiny\itshape]
  {alertes\\WebSocket} (6, -10.4);

\end{tikzpicture}

}
\caption{Architecture système en 5 couches — Smart Farm AI v3.0}
\label{fig:archi_systeme}
\end{figure}

\subsection{Description des couches}

\begin{description}
  \item[\textbf{Couche Clients :}] Applications utilisateur — Owner Dashboard (React SPA),
        Worker PWA (Progressive Web App avec Dexie.js pour l'offline), Map Center
        (Leaflet/OpenStreetMap) et Assistant (interface chat multimodale).

  \item[\textbf{Couche API Gateway :}] Backend FastAPI exposant 50+ endpoints REST,
        protégés par middleware JWT, avec support WebSocket pour les alertes temps réel
        et configuration CORS pour les domaines autorisés.

  \item[\textbf{Couche Services Métier :}] Logique applicative — CV Service (chargement
        et inférence YOLO), RAG Service (requêtage ChromaDB), LLM Agent (orchestration
        Ollama/Groq), Anomaly Service (IsolationForest) et Weather Service (OpenMeteo).

  \item[\textbf{Couche Données :}] Persistance multi-backend — PostgreSQL (production),
        SQLite (Oracle LITE\_MODE), ChromaDB (vecteurs RAG), MLflow Registry (modèles),
        DVC (données et pipelines) et Ollama (modèles LLM locaux).

  \item[\textbf{Couche IoT :}] Nœuds ESP32 envoyant des données télémétriques via
        HTTP POST toutes les 10 secondes. Simulés via Wokwi, opérationnels sur
        matériel réel sans modification du firmware.
\end{description}

\section{Architecture backend — Clean Architecture}

Le backend FastAPI implémente une variante de la \textbf{Clean Architecture}
\cite{fastapi2023}, garantissant la testabilité indépendante de chaque couche et
la substituabilité de l'infrastructure (PostgreSQL ↔ SQLite ↔ tests en mémoire).

\begin{figure}[H]
\centering
\resizebox{\textwidth}{!}{%
% ============================================================
%  Clean Architecture Backend — Smart Farm AI (FastAPI)
% ============================================================
\begin{tikzpicture}[
  x=1cm, y=1cm,
  layer/.style={
    rectangle, rounded corners=5pt,
    draw=#1!70!black, fill=#1!12,
    minimum width=13cm, minimum height=1.5cm,
    align=center, font=\small\bfseries\color{#1!80!black}
  },
  box/.style={
    rectangle, rounded corners=3pt,
    draw=#1!60!black, fill=#1!18,
    text width=2.4cm, minimum height=1.1cm,
    align=center, font=\tiny\bfseries
  },
  arrow/.style={->, thick, >=Stealth, #1},
  dbarrow/.style={->, thick, >=Stealth, dashed, #1}
]

% ─── Couche Routes / Endpoints ────────────────────────────────
\node[layer=farmblue, text width=13cm] (routes) at (0, 0) {};
\node[font=\tiny\color{farmblue!50}, above right=0pt] at (-6.3, 0.55)
  {Couche Présentation / API};
\node[box=farmblue] at (-4.5, 0) {Auth Routes\\JWT / OTP};
\node[box=farmblue] at (-1.5, 0) {CV Routes\\YOLO Scan};
\node[box=farmblue] at ( 1.5, 0) {IoT Routes\\Telemetry};
\node[box=farmblue] at ( 4.5, 0) {Bee Routes\\HiveIntel};

% ─── Couche Services ─────────────────────────────────────────
\draw[arrow=farmblue] (0, -0.75) -- (0, -1.55);
\node[layer=farmgreen, text width=13cm] (serv) at (0, -2.3) {};
\node[font=\tiny\color{farmgreen!50}, above right=0pt] at (-6.3, -1.55)
  {Couche Services (Logique Métier)};
\node[box=farmgreen] at (-4.5, -2.3) {Auth Service\\Hasher + OTP};
\node[box=farmgreen] at (-1.5, -2.3) {CV Service\\Inférence YOLO};
\node[box=farmgreen] at ( 1.5, -2.3) {IoT Service\\Anomaly + Alert};
\node[box=farmgreen] at ( 4.5, -2.3) {Bee Service\\Calculs apicoles};

% ─── Couche Repositories ─────────────────────────────────────
\draw[arrow=farmgreen] (0, -3.05) -- (0, -3.85);
\node[layer=farmorange, text width=13cm] (repo) at (0, -4.6) {};
\node[font=\tiny\color{farmorange!50}, above right=0pt] at (-6.3, -3.85)
  {Couche Repositories (Accès Données)};
\node[box=farmorange] at (-4.5, -4.6) {User Repo\\CRUD Users};
\node[box=farmorange] at (-1.5, -4.6) {Farm Repo\\CRUD Farms};
\node[box=farmorange] at ( 1.5, -4.6) {IoT Repo\\Telemetry CSV};
\node[box=farmorange] at ( 4.5, -4.6) {Bee Repo\\Hives + Visits};

% ─── Couche Models ───────────────────────────────────────────
\draw[arrow=farmorange] (0, -5.35) -- (0, -6.15);
\node[layer=farmred, text width=13cm] (models) at (0, -6.9) {};
\node[font=\tiny\color{farmred!50}, above right=0pt] at (-6.3, -6.15)
  {Couche Modèles (SQLAlchemy ORM)};
\node[box=farmred] at (-4.5, -6.9) {User\\Farm\\Worker};
\node[box=farmred] at (-1.5, -6.9) {AnimalUnit\\CVEvent};
\node[box=farmred] at ( 1.5, -6.9) {TelemetryRecord\\Alert};
\node[box=farmred] at ( 4.5, -6.9) {BeeHive\\BeeVisit\\BeeHarvest};

% ─── Couche Infrastructure ───────────────────────────────────
\draw[arrow=farmred] (0, -7.65) -- (0, -8.45);
\node[layer=farmgray, text width=13cm] (infra) at (0, -9.2) {};
\node[font=\tiny\color{farmgray!60}, above right=0pt] at (-6.3, -8.45)
  {Infrastructure (Bases de données \& Services externes)};
\node[box=farmgray] at (-5.5, -9.2) {PostgreSQL\\(prod)};
\node[box=farmgray] at (-2.8, -9.2) {SQLite\\(Oracle LITE)};
\node[box=farmgray] at ( 0,   -9.2) {ChromaDB\\(RAG)};
\node[box=farmgray] at ( 2.8, -9.2) {Ollama\\(LLM local)};
\node[box=farmgray] at ( 5.5, -9.2) {YOLO Models\\(.pt files)};

% ─── Annotations latérales ────────────────────────────────────
\draw[decorate, decoration={brace, amplitude=6pt, mirror},
      farmblue, thick] (7, -7.65) -- (7, -0.75)
  node[midway, right=8pt, font=\tiny\itshape\color{farmblue}]
  {Injection de dépendances\\(FastAPI Depends)};

\draw[decorate, decoration={brace, amplitude=6pt},
      farmgray, thick] (-7, -9.95) -- (-7, 0.75)
  node[midway, left=8pt, font=\tiny\itshape\color{farmgray}]
  {\rotatebox{90}{Indépendance des couches (DIP)}};

\end{tikzpicture}

}
\caption{Clean Architecture du backend FastAPI — Smart Farm AI v3.0}
\label{fig:clean_archi}
\end{figure}

\subsection{Organisation du code backend}

\begin{lstlisting}[style=pythonstyle, caption={Structure du projet backend FastAPI}, label={lst:backend_struct}]
backend/
├── app/
│   ├── api/v1/endpoints/
│   │   ├── auth_routes.py        # /auth/login, /auth/register, /auth/otp
│   │   ├── cv_routes.py          # /cv/scan, /cv/models, /cv/history
│   │   ├── iot_routes.py         # /iot/telemetry, /iot/latest, /iot/history
│   │   ├── bee_routes.py         # /bee/hives, /bee/visits, /bee/harvest
│   │   └── assistant_routes.py   # /assistant/chat, /assistant/tts
│   ├── core/
│   │   ├── config.py             # YOLO paths, Ollama URL, LITE_MODE flag
│   │   ├── security.py           # JWT encode/decode, OTP generation
│   │   └── database.py           # SQLAlchemy engine + session
│   ├── models/
│   │   └── domain.py             # 30+ SQLAlchemy models
│   ├── services/
│   │   ├── cv_service.py         # MODEL_REGISTRY, YOLO infer()
│   │   ├── rag_service.py        # ChromaDB query, embedding
│   │   └── anomaly_service.py    # IsolationForest score
│   └── repositories/
│       ├── user_repo.py
│       ├── farm_repo.py
│       └── telemetry_repo.py
├── mlops/
│   ├── dvc.yaml                  # Pipeline stages
│   └── scripts/
│       ├── build_dataset.py      # Stage 1 DVC
│       └── train_anomaly.py      # Stage 2 DVC → MLflow
└── requirements.txt
\end{lstlisting}

\section{Modèle de données}

\subsection{Entités principales}

Le modèle de données de \farmAI{} v3.0 comprend \textbf{30+ entités} organisées en
6 domaines fonctionnels.

\begin{table}[H]
\centering
\caption{Domaines fonctionnels et entités principales}
\label{tab:domains}
\begin{tabular}{lll}
\toprule
\textbf{Domaine} & \textbf{Entités} & \textbf{Description} \\
\midrule
Authentification & User, OTPToken, RefreshToken &
  Gestion des comptes, JWT, OTP email \\
Exploitations    & Farm, FarmField, WeatherCache &
  Fermes, parcelles, données météo \\
Télémétrie IoT   & TelemetryRecord, Alert, Node &
  Métriques ESP32, alertes, config nœuds \\
Vision IA        & CVEvent, CVModel, Species &
  Historique scans YOLO, modèles enregistrés \\
Apiculture       & BeeHive, BeeYard, BeeVisit,\newline
                   BeeHarvest, BeeStock, Expense &
  Gestion complète des ruchers \\
Ouvriers         & Worker, WorkerTask, WorkerGPS &
  Personnel, tâches, tracking GPS \\
\bottomrule
\end{tabular}
\end{table}

\subsection{Entités clés}

\begin{lstlisting}[style=pythonstyle, caption={Entités principales SQLAlchemy (extrait)}, label={lst:models}]
class User(Base):
    __tablename__ = "users"
    id         = Column(Integer, primary_key=True)
    email      = Column(String, unique=True, nullable=False)
    hashed_pw  = Column(String, nullable=False)
    role       = Column(Enum("owner","worker","admin"), default="owner")
    farms      = relationship("Farm", back_populates="owner")

class TelemetryRecord(Base):
    __tablename__ = "telemetry_records"
    id         = Column(Integer, primary_key=True)
    node       = Column(String, nullable=False)  # NODE_A / NODE_B
    metric     = Column(String, nullable=False)  # soil, weight, hive_temp...
    value      = Column(Float,  nullable=False)
    is_anomaly = Column(Boolean, default=False)
    timestamp  = Column(DateTime, default=datetime.utcnow)
    farm_id    = Column(Integer, ForeignKey("farms.id"))

class BeeHive(Base):
    __tablename__ = "bee_hives"
    id         = Column(Integer, primary_key=True)
    name       = Column(String, nullable=False)
    qr_code    = Column(String, unique=True)
    yard_id    = Column(Integer, ForeignKey("bee_yards.id"))
    visits     = relationship("BeeVisit", back_populates="hive")
\end{lstlisting}

\section{Architecture IoT}

\subsection{NODE\_A — Nœud d'irrigation}

Le NODE\_A surveille l'irrigation intelligente avec les capteurs et seuils suivants :

\begin{table}[H]
\centering
\caption{Configuration NODE\_A — Irrigation ESP32}
\label{tab:node_a}
\begin{tabular}{llll}
\toprule
\textbf{Pin GPIO} & \textbf{Capteur} & \textbf{Métrique} & \textbf{Seuils} \\
\midrule
34 (ADC)  & Capteur sol (résistif)  & \code{soil}      & Sec $<$35\% / Humide $>$65\% \\
35 (ADC)  & Capteur pression        & \code{pressure}  & Max : 10 bar \\
32 (ADC)  & Débitmètre YF-S201      & \code{flow}      & Fuite $<$0.5 L/min (pompe ON) \\
14 (DHT)  & DHT22 (température)     & \code{temp}      & Référence extérieure \\
27 (OUT)  & Relais pompe            & \code{pump}      & 1=ON / 0=OFF \\
26 (OUT)  & Electrovanne            & \code{valve}     & 1=Ouverte / 0=Fermée \\
\bottomrule
\end{tabular}
\end{table}

\subsection{NODE\_B — Nœud de surveillance des ruches}

\begin{table}[H]
\centering
\caption{Configuration NODE\_B — Rucher ESP32}
\label{tab:node_b}
\begin{tabular}{llll}
\toprule
\textbf{Pin GPIO} & \textbf{Capteur} & \textbf{Métrique} & \textbf{Seuils d'alerte} \\
\midrule
34 (ADC) & Capteur de poids (jauge) & \code{weight}    & SWARM si $\Delta >$2 kg/h \\
14 (DHT) & DHT22 extérieur          & \code{ext\_temp, ext\_hum} & Hum $>$85\% WARN \\
15 (1-W) & DS18B20 (ruche)          & \code{hive\_temp}& $<$30°C ou $>$38°C ALERT \\
2  (OUT) & LED rouge alerte         & --               & SWARM\_RISK / THERMAL \\
4  (OUT) & LED verte production     & --               & HONEY\_FLOW positif \\
\bottomrule
\end{tabular}
\end{table}

\subsection{Flux de données IoT}

\begin{lstlisting}[style=arducstyle, caption={Envoi de télémétrie HTTP POST (extrait NODE\_B)}, label={lst:iot_post}]
void postMetric(const char* metric, float value) {
    HTTPClient http;
    String url = String("http://") + API_HOST + ":"
               + API_PORT + "/api/v1/iot/telemetry";
    http.begin(url);
    http.addHeader("Content-Type", "application/json");
    http.setTimeout(3000);
    // JSON body : node + metric + value
    String body = "{\"node\":\"NODE_B\","
                  "\"metric\":\"" + String(metric) + "\","
                  "\"value\":" + String(value, 4) + "}";
    int code = http.POST(body);
    http.end();
}
\end{lstlisting}

Le flux complet est le suivant :
\[
\text{ESP32 (Wokwi)} \xrightarrow{\text{WiFi}} \text{HTTP POST}
\xrightarrow{} \text{FastAPI /iot/telemetry}
\xrightarrow{\text{SQLAlchemy}} \text{PostgreSQL/SQLite}
\xrightarrow{\text{IsolationForest}} \text{Anomaly Alert}
\xrightarrow{\text{WebSocket}} \text{Dashboard}
\]

\section{Architecture de déploiement}

\farmAI{} v3.0 supporte deux modes de déploiement :

\begin{table}[H]
\centering
\caption{Modes de déploiement}
\label{tab:deploy}
\begin{tabular}{lll}
\toprule
\textbf{Mode} & \textbf{Local (dev)} & \textbf{Oracle Cloud (prod)} \\
\midrule
Base de données  & PostgreSQL Docker   & SQLite (LITE\_MODE=true) \\
LLM              & Ollama Labess-7B    & Mock KB (sans GPU) \\
YOLO             & 8 modèles .pt       & Désactivé (cloud ARM64) \\
Frontend         & Vite dev server     & Nginx (Caddy HTTPS) \\
Orchestration    & \code{start\_all.ps1} & Docker Compose \\
SSL              & HTTP local          & Caddy ACME automatique \\
\bottomrule
\end{tabular}
\end{table}
